% concatenated from maximal_central_function_v01.tex
\documentclass[12pt]{amsart}



\headheight=8pt     \topmargin=0pt
\textheight=624pt \textwidth=432pt \oddsidemargin=18pt
\evensidemargin=18pt
\numberwithin{equation}{section}
\usepackage{graphicx}
\usepackage{amsfonts}
\usepackage{amssymb}
\usepackage{amsmath}
\usepackage{lmodern}
\usepackage{wrapfig}
\usepackage{placeins}
\usepackage{float}



\begin{document}
 
 \title{ On the Central Maximal Function in $\mathbb{R}^n$  }
 
\author{ A. Solyanik}

\date{\today}
\dedicatory{To my dear friend Alexander Stokolos}

\address{Odessa National Polytechnic University/Odessa/ Ukraine}
\email{solyanik@opu.ua}

\newtheorem{remark}{Remark}
\newtheorem{theorem}{Theorem}
\newtheorem{lemma}{Lemma}
\newtheorem{lemma*}{Lemma*}
\newtheorem{proposition}{Proposition}
\newtheorem{corollary}{Corollary}
\newtheorem{example}{Example}
\maketitle
\begin{abstract}
We study the central Hardy--Littlewood maximal operator $M_c$ in
$\mathbb{R}^n$ for all $n\geq1$, with particular emphasis on the question
of its injectivity. More precisely, we consider whether there exist two
nonnegative functions $f,g\in L^1(\mathbb{R}^n)$ such that
$\|f-g\|_1>0$ and $M_cf=M_cg$.
Along the way, we obtain an elementary explicit formula for $M_c(\chi_U)$,
where $U$ is the unit ball in $\mathbb{R}^3$.
We also derive an explicit formula, in terms of the incomplete beta
function, for the $n$-dimensional volume of the intersection of two balls
with arbitrary radii and centers. We include this elementary formula
because we were unable to find it in standard reference books.
\end{abstract}

\section{Introduction}
Let $B(x, r) = \{y \in \mathbb{R}^n : \lVert x - y \rVert < r\}$ be the ball
of radius $r$ centered at $x$. For $f\in L^1_\text{loc}(\mathbb{R}^n)$, the
central Hardy--Littlewood maximal function is defined by
\[
M_cf(x)=\sup_{r>0} \frac{1}{\vert B(x, r) \vert}\int_{B(x,r)} \vert f(t) \vert\,dt.
\]
For background on this operator and the differentiation of integrals,
see, for example, de~Guzm\'an~\cite{deguzman1975}.

During the author's visit to Baylor University in September 2023,
Professor Paul Hagelstein asked whether two nonnegative functions in
$L^1(\mathbb{R}^n)$ that differ on a set of positive measure can have the
same central maximal function. This paper provides a positive answer to
this question.

The uniqueness problem for maximal operators was studied by Lasha Ephremidze
in a series of papers; a survey of these results and further references can be
found in \cite{ephremidze2010}. In particular, he proved that a nonnegative
function in $L^1(\mathbb{R})$ is uniquely determined, up to equality almost
everywhere, by its one-sided Hardy--Littlewood maximal function
\cite[Theorem~3]{ephremidze2010}. This result concerns the one-sided operator,
rather than the central maximal operator $M_c$ considered here.

We begin with the one-dimensional case. The idea is to place a small bump
of the form $\epsilon\chi_I$, where $I$ is an interval and $\epsilon>0$ is
small, in the gap between two fixed ``towers''. The key point is that this
bump can be moved within the gap without changing the central maximal
function of the resulting function.

As Professor Hagelstein communicated to the author in a private letter
dated March~14, 2025, the same idea in the one-dimensional case was proposed
independently by Fedor Nazarov.

\section{The One-Dimensional Result}
\label{one_dimensional_result}

\begin{theorem}
\label{one_dimensional_theorem}
Let $0<\epsilon\leq1/4$ and $0<\alpha<\epsilon$. Set
\[
E=(-1-\epsilon,-\epsilon)\cup(\epsilon,1+\epsilon),
\qquad I=(\alpha-\epsilon,\alpha),
\]
and let
\[
f(x)=\chi_E(x)+\epsilon\chi_I(x).
\]
Then, with
\[
x_\epsilon=\frac{1+3\epsilon+\epsilon^3}{1+\epsilon^2},
\]
we have
\begin{equation}
\label{one_dimensional_maximal}
M_cf(x)=
\begin{cases}
\displaystyle 1-\frac{\epsilon-\epsilon^2/2}{1+\epsilon-|x|},
    & |x|\leq\epsilon,\\[8pt]
1,  & \epsilon<|x|<1+\epsilon,\\[4pt]
\displaystyle \frac{1}{2(|x|-\epsilon)},
    & 1+\epsilon\leq|x|\leq x_\epsilon,\\[8pt]
\displaystyle \frac{2+\epsilon^2}{2(|x|+1+\epsilon)},
    & |x|>x_\epsilon.
\end{cases}
\end{equation}
In particular, $M_cf$ does not depend on $\alpha$.
\end{theorem}

\begin{proof}
Write
\[
A(x,r)=\frac{1}{2r}\int_{x-r}^{x+r}f(t)\,dt.
\]
Reflection replaces $\alpha$ by $\epsilon-\alpha$, which satisfies the
same hypotheses. It therefore suffices to prove the formula for $x\geq0$,
uniformly in $\alpha$.

First, let $0\leq x\leq\epsilon$. If $0<r\leq\epsilon+x$, the
interval $(x-r,x+r)$ meets at most one tower, whose contribution has
length at most $r$. Since $f\leq\epsilon$ off the towers,
\[
A(x,r)\leq\frac{1+\epsilon}{2}
<1-\epsilon+\frac{\epsilon^2}{2}.
\]
For $\epsilon+x\leq r\leq1+\epsilon-x$, the interval contains $I$
and meets both towers without extending beyond their outer endpoints.
Hence
\[
A(x,r)=1-\frac{\epsilon-\epsilon^2/2}{r},
\]
which is increasing in $r$. For $1+\epsilon-x\leq r\leq1+\epsilon+x$,
\[
A(x,r)=\frac12+\frac{1-\epsilon+\epsilon^2-x}{2r},
\]
which is decreasing, since $1-\epsilon+\epsilon^2-x>0$.
For $r\geq1+\epsilon+x$, the average is $(2+\epsilon^2)/(2r)$ and
is again decreasing. Thus the maximum is attained at $r=1+\epsilon-x$
and equals the first expression in \eqref{one_dimensional_maximal}.

For $\epsilon<x<1+\epsilon$, we have $M_cf(x)=1$, since $f=1$ in a
neighborhood of $x$ and $0\leq f\leq1$ everywhere.

It remains to consider $x\geq1+\epsilon$. Set
\[
r_1=x-\epsilon,\qquad r_2=x+1+\epsilon.
\]
At these radii the interval contains, respectively, the entire right
tower alone and the entire support of $f$. Consequently,
\[
A(x,r_1)=\frac{1}{2(x-\epsilon)},\qquad
A(x,r_2)=\frac{2+\epsilon^2}{2(x+1+\epsilon)}.
\]
For $0<r\leq r_1$, only the right tower can be encountered and its
average is nondecreasing. For $r\geq r_2$, the total mass is fixed and
the average is decreasing. To handle the intermediate radii, put
\[
N(r)=\int_{x-r}^{x+r}f(t)\,dt,
\qquad k=\frac{1+\epsilon^2}{1+2\epsilon}.
\]
Our assumption on $\epsilon$ implies $\epsilon\leq k<1$.
For $r_1\leq r\leq x+\epsilon$, the right tower is fully contained,
the left tower is absent, and the remaining density is at most $\epsilon$.
Thus
\[
N(r)\leq1+\epsilon(r-r_1)\leq1+k(r-r_1).
\]
For $x+\epsilon\leq r\leq r_2$, we have
\[
N(r)=2+\epsilon^2-(r_2-r)
\leq2+\epsilon^2-k(r_2-r)=1+k(r-r_1).
\]
Therefore, throughout $[r_1,r_2]$,
\[
A(x,r)\leq\frac{k}{2}+\frac{1-kr_1}{2r}.
\]
The right-hand side is monotone or constant in $r$ and agrees with
$A(x,r)$ at both endpoints. It follows that
\begin{equation}
\label{one_dimensional_maximal_exterior}
M_cf(x)=\max\left\{
\frac{1}{2(x-\epsilon)},
\frac{2+\epsilon^2}{2(x+1+\epsilon)}
\right\},\qquad x\geq1+\epsilon.
\end{equation}
The two expressions are equal precisely when $x=x_\epsilon$; the first
is larger for $x<x_\epsilon$ and the second for $x>x_\epsilon$.
Since
\[
x_\epsilon-(1+\epsilon)
=\frac{\epsilon(2-\epsilon)}{1+\epsilon^2}>0,
\]
this proves the remaining cases of \eqref{one_dimensional_maximal}.
\end{proof}

For example, if $0<\alpha<\beta<\epsilon$ and
\[
g(x)=\chi_E(x)+\epsilon\chi_{(\beta-\epsilon,\beta)}(x),
\]
then $f,g\in L^1(\mathbb{R})$ are nonnegative and
\[
M_cf(x)=M_cg(x)\quad\text{for every }x\in\mathbb{R},
\qquad \|f-g\|_1=2\epsilon(\beta-\alpha)>0.
\]

\section{The Volume of the Hyperspherical Cap}

For the subsequent analysis, we need convenient formulas for the volumes
of a hyperspherical cap and the intersection of two balls in
$\mathbb{R}^n$. Since we were unable to find the corresponding formulas in
authoritative reference works, we present them here, together with
complete proofs.

In this and the next section, we assume that $n\geq2$, since our subsequent
analysis concerns higher-dimensional problems. The corresponding formulas
remain valid for $n=1$ and are trivial to verify in that case.

We use the following notation and definitions.
By $E^c$ we denote the set $\lbrace x\in\mathbb{R}^n : x\notin E\rbrace$,
and by $d(x,E)=\inf\lbrace\Vert x-y\Vert : y\in E\rbrace$ we denote the
distance from the point $x$ to the set $E$.

The incomplete beta function is defined by
\begin{equation}
\label{incomplete_beta}
\mathrm{B}(a,b;z)=\int_0^z t^{a-1}(1-t)^{b-1}\,dt,
\qquad a,b>0,\quad 0\leq z\leq1.
\end{equation}
The complete beta function is obtained when $z=1$ and can be expressed
in terms of the gamma function as
\begin{equation}
\label{complete_beta}
\mathrm{B}(a,b)=\mathrm{B}(a,b;1)
=\frac{\Gamma(a)\Gamma(b)}{\Gamma(a+b)},
\end{equation}
where
\[
\Gamma(z)=\int_0^\infty t^{z-1}e^{-t}\,dt,\qquad z>0.
\]

The volume of the $n$-dimensional ball $B_r=B(x,r)$ is given by
\begin{equation}
\label{ball_volume}
|B_r|=c_n r^n,
\end{equation}
where, here and throughout the paper, $|E|$ denotes the $n$-dimensional
Lebesgue measure of a measurable set $E$, and
\begin{equation}
c_n=\frac{\pi^{n/2}}{\Gamma\!\left(\frac n2+1\right)}.
\end{equation}

\begin{remark}
In the Python library \texttt{scipy.special}, the function
\begin{equation}
\label{python_beta}
\texttt{betainc}(a,b,z)
=\frac{\mathrm{B}(a,b;z)}{\mathrm{B}(a,b)}
\end{equation}
returns the regularized incomplete beta function. Thus its value must be
multiplied by $\mathrm{B}(a,b)$ to obtain the function
$\mathrm{B}(a,b;z)$ used here.
\end{remark}

Let $H$ be an open half-space in $\mathbb{R}^n$, written as
\[
H=z+\{y\in\mathbb{R}^n:\langle y,e\rangle>0\},
\qquad z\in\mathbb{R}^n,\quad \|e\|=1.
\]
Here $\langle x,y\rangle=\sum_{j=1}^n x_jy_j$ is the scalar product.
The boundary
\[
\partial H=z+\{y\in\mathbb{R}^n:\langle y,e\rangle=0\}
\]
is an $(n-1)$-dimensional hyperplane.

For a ball $B=B(x,r)$, its intersection $C_h=B\cap H$ is called a
\emph{hyperspherical cap}. We denote by $h$ the \emph{signed distance}
from the center $x$ to the cutting hyperplane $\partial H$, with the
following sign convention:
\[
h=
\begin{cases}
d(x,\partial H),&x\notin H,\\
-d(x,\partial H),&x\in H.
\end{cases}
\]
Thus $h<0$ when the center lies in the half-space defining the cap,
and $h=0$ when the cutting hyperplane passes through the center.
For $-r\leq h\leq r$, the usual geometric height of the cap is $r-h$.

By translation and rotation invariance, the volume of a cap depends only
on $r$ and $h$. We may therefore take $B_r=B(0,r)$ and
$H=\{x\in\mathbb{R}^n:x_n>h\}$.

\begin{figure}[htbp]
 \centering
 \includegraphics[width=0.8\textwidth]{cap.jpg}
 \caption{Hyperspherical cap $C_r^h$, with signed distance $h$ and height $r-h$.}
 \label{fig:spherical_cap}
\end{figure}

\begin{lemma}
Let $r>0$, $-r\leq h\leq r$, and
\[
C_r^h=\{x=(x_1,\ldots,x_n)\in B_r:x_n>h\},
\qquad B_r=B(0,r).
\]
Then
\begin{equation}
\label{cap_volume}
|C_r^h|=c_{n-1}(2r)^n\,
\mathrm{B}\!\left(
\frac{n+1}{2},\frac{n+1}{2};
\frac12\left(1-\frac hr\right)\right).
\end{equation}
In particular,
\[
|C_r^r|=0,\qquad |C_r^{-r}|=c_n r^n,
\qquad |C_r^0|=\frac12c_n r^n.
\]
\end{lemma}

\begin{proof}
The section of $B_r$ by the hyperplane $x_n=t$, where $|t|<r$,
is an $(n-1)$-dimensional ball of radius $\sqrt{r^2-t^2}$.
Integrating these sections and then making the substitution
$t=r(1-2v)$, we obtain
\[
\begin{aligned}
|C_r^h|
&=c_{n-1}\int_h^r(r^2-t^2)^{(n-1)/2}\,dt\\
&=c_{n-1}(2r)^n
  \int_0^{(1-h/r)/2}v^{(n-1)/2}(1-v)^{(n-1)/2}\,dv.
\end{aligned}
\]
The last integral is the incomplete beta function in
\eqref{cap_volume}. The three particular values also follow directly
from the geometry of the corresponding caps.
\end{proof}

For instance, when $n=3$, evaluating $\mathrm{B}(2,2;z)$ gives
\begin{equation}
\label{cap_volume_three}
\begin{aligned}
\mathrm{Vol}(C_r^h)
&=8\pi r^3\mathrm{B}\!\left(2,2;\frac12\left(1-\frac hr\right)\right)\\
&=\frac{2\pi r^3}{3}-\pi r^2h+\frac{\pi h^3}{3}
 =\frac{\pi}{3}(r-h)^2(2r+h).
\end{aligned}
\end{equation}
\FloatBarrier

\section{The Volume of the Intersection of Two Balls}

Consider two balls $B_1=B(x_1,r_1)$ and $B_2=B(x_2,r_2)$ in
$\mathbb{R}^n$, with $r_1,r_2>0$ and $d=\|x_1-x_2\|$.
If $d\leq|r_1-r_2|$, one ball is contained in the other. If
$d\geq r_1+r_2$, the open balls are disjoint. Thus
\[
|B_1\cap B_2|=
\begin{cases}
c_n\min\{r_1,r_2\}^n,&d\leq|r_1-r_2|,\\[2pt]
0,&d\geq r_1+r_2.
\end{cases}
\]
The first case includes coincident centers, $d=0$. It remains to treat
partial overlap, for which the following formula applies.

\begin{theorem}
Let $B_1$ and $B_2$ be two balls in $\mathbb{R}^n$ with positive radii
$r_1$ and $r_2$ and centers at distance $d$ apart, where
\begin{equation}
\label{condition}
|r_1-r_2|<d<r_1+r_2.
\end{equation}
Then
\begin{equation}
\label{vol_inter_two_balls}
\begin{aligned}
|B_1\cap B_2|
={}&\frac{2^n\pi^{(n-1)/2}}{\Gamma((n+1)/2)}
\Bigg[
 r_1^n\,\mathrm{B}\!\left(
 \frac{n+1}{2},\frac{n+1}{2};
 \frac12\left(1-\frac{h_1}{r_1}\right)\right)\\
&\hspace{55pt}{}+
 r_2^n\,\mathrm{B}\!\left(
 \frac{n+1}{2},\frac{n+1}{2};
 \frac12\left(1-\frac{h_2}{r_2}\right)\right)
\Bigg],
\end{aligned}
\end{equation}
where
\begin{align}
h_1&=\frac{r_1^2-r_2^2+d^2}{2d},\label{h_mu}\\
h_2&=\frac{r_2^2-r_1^2+d^2}{2d}.\label{h_r}
\end{align}
These signed distances satisfy $-r_j<h_j<r_j$ for $j=1,2$, so both
arguments of the incomplete beta function lie in $(0,1)$.
\end{theorem}

\begin{proof}
Write $B_k=B(x_k,r_k)$ for $k=1,2$. Without loss of generality,
assume that $r_2\geq r_1$.
Under \eqref{condition}, the common boundary
$\partial B_1\cap\partial B_2$ is a nonempty $(n-2)$-dimensional
sphere. Let $L$ be the hyperplane containing it; this hyperplane is
perpendicular to the line through $x_1$ and $x_2$.
Let $H_1$ be the open half-space bounded by $L$ that contains $x_2$,
and let $H_2$ be the opposite open half-space.

Let $h_1$ and $h_2$ be the signed distances from $x_1$ and $x_2$ to
$L$, with the signs determined by $H_1$ and $H_2$, respectively.
Then
\begin{equation}
\label{sum}
h_1+h_2=d.
\end{equation}
By the Pythagorean theorem, the squared radius of the common
$(n-2)$-dimensional sphere is both $r_1^2-h_1^2$ and $r_2^2-h_2^2$.
Consequently,
\[
r_1^2-h_1^2=r_2^2-(d-h_1)^2,
\]
which gives
\begin{equation}
\label{h_1}
2dh_1=d^2+r_1^2-r_2^2.
\end{equation}
Combining this with \eqref{sum}, we obtain
\begin{equation}
\label{h_2}
2dh_2=d^2+r_2^2-r_1^2.
\end{equation}
Thus \eqref{h_mu} and \eqref{h_r} follow. Since the common sphere has
positive radius, $|h_j|<r_j$ for $j=1,2$.

Set $C_{h_1}=B_1\cap H_1$ and $C_{h_2}=B_2\cap H_2$.
To verify that these caps make up the overlap, put
$e=(x_2-x_1)/d$. Then
\[
\begin{aligned}
&\bigl(\|y-x_2\|^2-r_2^2\bigr)
 -\bigl(\|y-x_1\|^2-r_1^2\bigr)\\
&\qquad=2d\bigl(h_1-\langle y-x_1,e\rangle\bigr).
\end{aligned}
\]
The right-hand side is negative on $H_1$ and positive on $H_2$.
It follows that $C_{h_1}=B_1\cap B_2\cap H_1$ and
$C_{h_2}=B_1\cap B_2\cap H_2$. These two caps are disjoint, and
the portion of the overlap in $L$ has zero $n$-dimensional measure.
Hence
\[
|B_1\cap B_2|=|C_{h_1}|+|C_{h_2}|.
\]
Applying \eqref{cap_volume} to the two caps and adding their volumes
yields \eqref{vol_inter_two_balls}.
\end{proof}

\begin{figure}[htbp]
 \centering
 \includegraphics[width=0.8\textwidth]{plot_balls.jpg}
 \caption{Intersection of the balls $B_1$ and $B_2$.}
 \label{fig:two_balls_intersection}
\end{figure}

\begin{corollary}[$n=2$]
For two disks $D_1$ and $D_2$ with radii $r_1,r_2$ and centers at
distance $d$ apart, where $|r_1-r_2|<d<r_1+r_2$, we have
\begin{equation}
\begin{aligned}
\text{Area}(D_1\cap D_2)
={}&r_1^2\arccos\!\left(\frac{d^2+r_1^2-r_2^2}{2dr_1}\right)\\
&{}+r_2^2\arccos\!\left(\frac{d^2+r_2^2-r_1^2}{2dr_2}\right)\\
&{}-\frac12\sqrt{(r_1+r_2-d)(r_1-r_2+d)
                 (r_2-r_1+d)(r_1+r_2+d)}.
\end{aligned}
\end{equation}
\end{corollary}

\begin{proof}
In dimension two, integration of the sections of a cap gives
\[
|C_r^h|=2\int_h^r\sqrt{r^2-t^2}\,dt
=r^2\arccos(h/r)-h\sqrt{r^2-h^2}.
\]
Add this expression for the two caps and use
$h_1+h_2=d$ and $r_1^2-h_1^2=r_2^2-h_2^2$.
The square-root contribution is $-d\sqrt{r_1^2-h_1^2}$,
and the claimed expression follows from the identity
\begin{align*}
4d^2(r_1^2-h_1^2)
={}&(r_1+r_2-d)(r_1-r_2+d)\\
&\qquad{}\times(r_2-r_1+d)(r_1+r_2+d).\qedhere
\end{align*}
\end{proof}

\begin{corollary}[$n=3$]
For two balls $B_1$ and $B_2$ with radii $r_1,r_2$ and centers at
distance $d$ apart, where $|r_1-r_2|<d<r_1+r_2$, we have
\begin{equation}
\label{vol_3_balls_inter}
\begin{aligned}
\text{Vol}(B_1\cap B_2)
={}&\frac{\pi}{12d}(r_1+r_2-d)^2\\
&\quad{}\times\bigl(d^2+2d(r_1+r_2)-3(r_1-r_2)^2\bigr).
\end{aligned}
\end{equation}
\end{corollary}

\begin{proof}
By \eqref{cap_volume_three}, the sum of the two cap volumes is
\[
\text{Vol}(B_1\cap B_2)
=\frac{\pi}{3}\sum_{j=1}^2(r_j-h_j)^2(2r_j+h_j).
\]
Substituting \eqref{h_mu} and \eqref{h_r} and simplifying gives
\eqref{vol_3_balls_inter}.
\end{proof}
\FloatBarrier


\section{The Two-Dimensional Result}
\label{two_dimensional_result}

We now implement the same idea in the plane, using an annulus in place
of the two towers. Throughout this section, let
\[
\delta=\frac12,\qquad U=B(0,1),\qquad B_\delta=B(0,\delta),
\qquad U_\delta=\{y\in\mathbb{R}^2:\delta<\|y\|<1\},
\]
and write
\begin{equation}
\label{two_dimensional_average}
A(x,r)=\frac{|U_\delta\cap B(x,r)|}{\pi r^2},\qquad r>0.
\end{equation}
The average is radial in its center. In the scalar computations below,
we therefore also use $x\geq0$ to denote the distance of the center
from the origin, as in the later sections of this paper.

Figure~\ref{fig:my_plot} illustrates the annulus and an averaging disk.

\begin{figure}[htbp]
 \centering
 \includegraphics[width=0.8\textwidth]{hollow_ball.jpg}
 \caption{The annulus $U_\delta$, its central hole $B_\delta$, and an
 averaging disk $B(x,r)$. The shaded region is $U_\delta\cap B(x,r)$.
 This is a schematic diagram, not drawn to scale.}
 \label{fig:my_plot}
\end{figure}

The small added disk will be allowed to move in a fixed neighborhood
of the origin contained in the hole. The essential point is a gap in
the possible maximizing radii: a disk contributing to the maximal
function can be chosen either to miss this neighborhood or to contain
it completely. We do not need to determine all local maxima of
$A(x,\cdot)$, or to prove that there are exactly two of them.

The numerical profiles in Figure~\ref{fig:two_dimensional_jump}
illustrate this change in the maximizing radius: at $x=1.6$ the
smaller-radius peak is higher, whereas at $x=1.85$ the larger-radius
peak is higher. The uniform gap needed for the construction is proved
below without relying on the numerical plots.

\begin{figure}[p]
 \centering
 \includegraphics[width=0.96\textwidth]{average_two_dim_delta_0,5_x_1,6.png}
 \par\smallskip
 {\small (a) $x=1.6$.\par}
 \medskip
 \includegraphics[width=0.96\textwidth]{average_two_dim_delta_0,5_x_1,85.png}
 \par\smallskip
 {\small (b) $x=1.85$.\par}
 \caption{Graphs of $r\mapsto A(x,r)$ for $\delta=1/2$, illustrating
 the change from the smaller-radius maximum in (a) to the larger-radius
 maximum in (b). The dashed lines mark $r=x-1$, $r=x-\delta$,
 $r=x+\delta$, and $r=x+1$.}
 \label{fig:two_dimensional_jump}
\end{figure}
\FloatBarrier

\begin{theorem}
\label{two_dimensional_theorem}
There exist constants $0<\eta<\delta$ and $\epsilon_0>0$ such that,
for every $0<\epsilon\leq\epsilon_0$ and $0<\sigma<\eta$, the central
maximal function of
\[
f_a(y)=\chi_{U_\delta}(y)+\epsilon\chi_{B(a,\sigma)}(y)
\]
is independent of $a$ whenever $\|a\|+\sigma<\eta$.
In particular, if
\[
\sigma=\frac{\eta}{8},\qquad
 a=\left(\frac{\eta}{4},0\right),\qquad
 b=\left(-\frac{\eta}{4},0\right),
\]
then $f=f_a$ and $g=f_b$ are bounded, nonnegative, compactly supported
functions in $L^1(\mathbb{R}^2)$ satisfying
\[
M_cf(x)=M_cg(x)\quad\text{for every }x\in\mathbb{R}^2,
\qquad \|f-g\|_1=2\epsilon\pi\sigma^2>0.
\]
\end{theorem}

We first exclude the radius $r=x$, and then obtain a uniform exclusion
interval which persists after a small perturbation.

\begin{lemma}
\label{two_dimensional_diagonal_gap}
For every $x>0$, we have
\begin{equation}
\label{two_dimensional_strict_gap}
A(x,x)<M_c(\chi_{U_\delta})(x).
\end{equation}
\end{lemma}

\begin{proof}
For $s>0$, put $J_s(x,r)=|B(0,s)\cap B((x,0),r)|$ and
\[
H_s(x)=x\frac{\partial J_s}{\partial r}(x,x)-2J_s(x,x).
\]
If $2x\leq s$, then $H_s(x)=0$. If $2x>s$, the
plane-disk intersection formula in the preceding section gives
\begin{align*}
J_s(x,x)
&=s^2\arccos\!\left(\frac{s}{2x}\right)
 +2x^2\arcsin\!\left(\frac{s}{2x}\right)
 -\frac{s}{2}\sqrt{4x^2-s^2},\\
\frac{\partial J_s}{\partial r}(x,x)
&=4x\arcsin\!\left(\frac{s}{2x}\right).
\end{align*}
Consequently,
\begin{equation}
\label{two_dimensional_Hs}
H_s(x)=s\sqrt{4x^2-s^2}
       -2s^2\arccos\!\left(\frac{s}{2x}\right),\qquad 2x>s.
\end{equation}
The first derivative with respect to $r$ is continuous at tangency
for $x>0$, so the preceding statements also cover $2x=s$.
Since the annulus average is $(J_1-J_\delta)/(\pi r^2)$,
\begin{equation}
\label{two_dimensional_diagonal_derivative}
\pi x^3\frac{\partial A}{\partial r}(x,x)=H_1(x)-H_\delta(x).
\end{equation}

If $0<x\leq1/4$, then $B((x,0),x)$ lies in the hole, and
$A(x,x)=0<M_c(\chi_{U_\delta})(x)$.

If $1/4<x\leq1/2$, then $H_1(x)=0$. Set
$\theta=\arccos(\delta/(2x))$, so that $0<\theta\leq\pi/3$.
Formula \eqref{two_dimensional_Hs} becomes
\[
H_\delta(x)=\delta^2(\tan\theta-2\theta)<0.
\]
Indeed, $\cos\theta\geq1/2$ and $\sin\theta<\theta$ imply
$\tan\theta\leq2\sin\theta<2\theta$.
Thus $\partial A/\partial r(x,x)>0$, which proves
\eqref{two_dimensional_strict_gap} in this case.

If $1/2<x<1$, the center lies in the annulus, and
$M_c(\chi_{U_\delta})(x)=1$. The disk $B((x,0),x)$ intersects the
hole in positive area, so $A(x,x)<1$.

Next, suppose that $1\leq x\leq2$, and put
$D(x)=H_1(x)-H_\delta(x)$. Define
\[
q(t)=\frac{1-2t}{\sqrt{1-t}}.
\]
Differentiating \eqref{two_dimensional_Hs}, and using $\delta=1/2$,
we obtain
\[
D'(x)=2q\!\left(\frac{1}{4x^2}\right)
          -q\!\left(\frac{1}{16x^2}\right).
\]
The function $q$ is decreasing on $[0,1/4]$, with $q(0)=1$ and
$q(1/4)=1/\sqrt3$. Hence
\[
D'(x)\geq\frac{2}{\sqrt3}-1>0\qquad(x\geq1).
\]
On the other hand,
\begin{align*}
D(2)
&=\sqrt{15}-\frac{\sqrt{63}}4-\frac{3\pi}{4}
  +2\arcsin\!\left(\frac14\right)
  -\frac12\arcsin\!\left(\frac18\right)\\
&\leq\frac{17}{\sqrt{15}}-\frac{\sqrt{63}}4
      -\frac{3\pi}{4}-\frac1{16}<0.
\end{align*}
Here we used $t\leq\arcsin t\leq t/\sqrt{1-t^2}$.
For completeness, the last strict inequality follows entirely from
rational bounds: $\sqrt{15}>58/15$, $\sqrt{63}>198/25$, and
$\pi>157/50$ give
\[
\frac{17}{\sqrt{15}}-\frac{\sqrt{63}}4-\frac{3\pi}{4}-\frac1{16}
<\frac{255}{58}-\frac{99}{50}-\frac{471}{200}-\frac1{16}
=-\frac{11}{11600}.
\]
Thus $D(x)\leq D(2)<0$ for $1\leq x\leq2$. By
\eqref{two_dimensional_diagonal_derivative}, decreasing the radius
slightly increases the average, and \eqref{two_dimensional_strict_gap}
follows.

Finally, let $x\geq2$. Integration in polar coordinates gives
\begin{align*}
\pi x^2 A(x,x)
&=2\int_{1/2}^1 s\arccos\!\left(\frac{s}{2x}\right)\,ds\\
&=\frac{3\pi}{8}
  -2\int_{1/2}^1 s\arcsin\!\left(\frac{s}{2x}\right)\,ds\\
&\leq\frac{3\pi}{8}-\frac{7}{24x}.
\end{align*}
The disk of radius $x+1$ contains the whole annulus, so
\[
A(x,x+1)=\frac{3}{4(x+1)^2}.
\]
This is strictly larger than the preceding upper bound for $A(x,x)$.
Indeed,
\[
\frac{3}{4(x+1)^2}
-\left(\frac{3}{8x^2}-\frac{7}{24\pi x^3}\right)
=\frac{P(x)}{24\pi x^3(x+1)^2},
\]
where
\begin{align*}
P(x)={}&9\pi(x-2)^3+(36\pi+7)(x-2)^2\\
       &{}+(27\pi+42)(x-2)+63-18\pi>0
       \qquad(x\geq2).
\end{align*}
All coefficients in this last expression are positive. This completes
the proof of \eqref{two_dimensional_strict_gap}.
\end{proof}

\begin{lemma}
\label{two_dimensional_uniform_gap}
There exist $0<\eta\leq1/16$ and $0<\epsilon_0\leq1/4$ with the
following property. If $p$ is measurable, $0\leq p\leq\epsilon_0$
almost everywhere, and $p=0$ almost everywhere outside $B_\eta=B(0,\eta)$,
then the central maximal function of $\chi_{U_\delta}+p$ can be computed
using only radii satisfying
\begin{equation}
\label{two_dimensional_safe_radii}
\bigl|r-\|x\|\bigr|>\eta.
\end{equation}
More precisely, for each center $x$, the supremum over the other
positive radii is strictly smaller than the supremum over the radii
in \eqref{two_dimensional_safe_radii}.
\end{lemma}

\begin{proof}
First consider centers with $1/8\leq x\leq10$, where $x$ denotes the
distance from the origin. For each fixed $r>0$, $A(x,r)$ is continuous
in $x$. Thus $M_c(\chi_{U_\delta})=\sup_{r>0}A(\cdot,r)$ is lower
semicontinuous. The function $A(x,x)$ is continuous for $x>0$.
Lemma~\ref{two_dimensional_diagonal_gap} therefore implies that
\begin{equation}
\label{two_dimensional_gamma}
\gamma=\min_{1/8\leq x\leq10}
\bigl(M_c(\chi_{U_\delta})(x)-A(x,x)\bigr)>0.
\end{equation}
Also $\gamma\leq1$, since $0\leq A\leq M_c(\chi_{U_\delta})\leq1$.
Choose
\begin{equation}
\label{two_dimensional_constants}
\eta=\frac{\gamma}{128},\qquad
\epsilon_0=\frac{\gamma}{4}.
\end{equation}
In particular, $\eta\leq1/16$ and $\epsilon_0\leq1/4$.

Since $0\leq\chi_{U_\delta}\leq1$, differentiation of its disk average
gives the bound
\[
\left|\frac{\partial A}{\partial r}(x,r)\right|\leq\frac2r.
\]
This also follows directly by writing the derivative as the difference
of the boundary and interior contributions, each of which lies between
$0$ and $2/r$. For $1/8\leq x\leq10$ and $|r-x|\leq\eta$, all radii
between $x$ and $r$ are at least $1/16$. Hence
\[
|A(x,r)-A(x,x)|\leq32\eta=\frac{\gamma}{4}.
\]
It follows that every such average of the perturbed function satisfies
\begin{align*}
\frac1{\pi r^2}\int_{B(x,r)}(\chi_{U_\delta}+p)
&\leq A(x,x)+\frac{\gamma}{4}+\epsilon_0\\
&\leq M_c(\chi_{U_\delta})(x)-\frac{\gamma}{2}.
\end{align*}
For $p=0$, the same estimate without the term $\epsilon_0$ shows that
the supremum defining $M_c(\chi_{U_\delta})(x)$ can already be taken
over $|r-x|>\eta$. Since $p\geq0$, the supremum of the perturbed
averages over these radii is at least $M_c(\chi_{U_\delta})(x)$.
This proves the required strict gap for the compact range of centers.

For $0\leq x\leq1/8$ and $|r-x|\leq\eta$, we have
\[
x+r\leq2x+\eta\leq\frac5{16}<\delta.
\]
Thus the averaging disk is contained in the hole and its average is
at most $\epsilon_0\leq1/4$. The safe radius $x+1$ gives an average
at least
\[
\frac{3}{4(x+1)^2}\geq\frac{16}{27}>\frac14.
\]

It remains to treat $x\geq10$. Rotate coordinates so that the center
is $(x,0)$, and write $r=x+t$, where $|t|\leq\eta$.
Every $y=(y_1,y_2)\in B((x,0),r)$ satisfies
\[
y_1>\frac{\|y\|^2-2xt-t^2}{2x}
\geq-\eta-\frac{\eta^2}{2x}.
\]
The part of the annulus with $y_1>0$ has area $3\pi/8$.
The additional strip has width at most $\eta+\eta^2/(2x)$, and its
sections in the unit disk have length at most $2$. Consequently,
\begin{align*}
\int_{B((x,0),r)}(\chi_{U_\delta}+p)
&\leq\frac{3\pi}{8}+2\eta+\frac{\eta^2}{x}
       +\epsilon_0\pi\eta^2\\
&\leq\frac{3\pi}{8}+\frac18+\frac1{2560}+\frac{\pi}{1024}
 <\frac{\pi}{2}.
\end{align*}
Therefore the average at an excluded radius is smaller than
\[
\frac{1}{2(x-\eta)^2}
<\frac{3}{4(x+1)^2}\qquad(x\geq10,\ \eta\leq1/16).
\]
The expression on the right is a lower bound for the average at the
safe radius $x+1$, which contains the whole annulus and the support
of $p$. This finishes the proof for all centers.
\end{proof}

\begin{proof}[Proof of Theorem~\ref{two_dimensional_theorem}]
Use the constants in Lemma~\ref{two_dimensional_uniform_gap}.
If a radius satisfies \eqref{two_dimensional_safe_radii}, then its
disk either misses $B_\eta$ or contains $B_\eta$ entirely. Accordingly,
for any $p$ as in that lemma, with $m=\int_{\mathbb{R}^2}p$, we have
\[
\frac1{\pi r^2}\int_{B(x,r)}p=
\begin{cases}
0,&0<r<\|x\|-\eta,\\[2pt]
\displaystyle\frac{m}{\pi r^2},&r>\|x\|+\eta.
\end{cases}
\]
Taking the supremum over these radii gives the exact identity
\begin{equation}
\label{two_dimensional_common_maximal}
\begin{aligned}
M_c(\chi_{U_\delta}+p)(x)
=\max\Bigg\{&\sup_{0<r<\|x\|-\eta} A(x,r),\\
&\sup_{r>\|x\|+\eta}
\left(A(x,r)+\frac{m}{\pi r^2}\right)\Bigg\},
\end{aligned}
\end{equation}
where the first supremum is taken to be $0$ when its set of radii is
empty. The right-hand side depends on $p$ only through its total mass.

For $p=\epsilon\chi_{B(a,\sigma)}$, this mass is
$m=\epsilon\pi\sigma^2$, independent of the center $a$.
The hypothesis $\|a\|+\sigma<\eta$ places the whole added disk in
$B_\eta$, and Lemma~\ref{two_dimensional_uniform_gap} applies.
This proves the asserted pointwise equality of the maximal functions.
For the two particular centers in the theorem, the added disks are
disjoint, so
\[
\|f_a-f_b\|_1
=\epsilon\bigl(|B(a,\sigma)|+|B(b,\sigma)|\bigr)
=2\epsilon\pi\sigma^2>0.
\]
Both functions are bounded and supported in the unit disk, completing
the proof.
\end{proof}

\begin{remark}
Formula \eqref{two_dimensional_common_maximal} proves more than
translation invariance for a disk: any two nonnegative perturbations
supported in $B_\eta$, bounded by $\epsilon_0$, and having the same
total mass give exactly the same maximal function. In particular,
that maximal function is radial even when the perturbation is not.
The allowed region is the smaller disk $B_\eta$, not the entire hole
$B_\delta$.
\end{remark}

\begin{remark}
The maximal function is unchanged when the bump moves, but it is not
equal to that of the annulus alone. Indeed, at the origin,
\[
M_c(\chi_{U_\delta})(0)=\frac34,
\qquad M_cf_a(0)\geq\frac34+\epsilon\sigma^2>\frac34,
\]
where the lower bound for $M_cf_a(0)$ is obtained with $r=1$.
Thus the assertion concerns the position of the added mass, not its removal.
\end{remark}


\section{The Three-Dimensional Result}
\label{three_dimensional_result}

The planar construction does not extend to $\mathbb{R}^3$ simply by
replacing the annulus with a uniform spherical shell. This is an actual
obstruction, rather than a difficulty in proving the required gap:
for every such shell there is an exterior center at which the unique
maximizing sphere passes through the origin. Moving a small added ball
from one side of the origin to the other then changes the maximal
function. We first establish this fact, and then replace the single
shell by three thin concentric layers.

The constructions in this and the following section were proposed by
the language model during the interaction described in the
acknowledgements below.

As in the preceding section, in radial computations we use $x\geq0$
for the distance of the center from the origin and take the center to
be $xe_1$, where $e_1=(1,0,0)$. For a nonnegative function $W$, write
\[
A_W(x,r)=\frac{3}{4\pi r^3}\int_{B(xe_1,r)}W(y)\,dy.
\]
For a fixed background we also write $A(x,r)$, as before.

\subsection{Why one spherical shell is insufficient}

Let $0<\delta<1$, let $U=B(0,1)$ and $B_\delta=B(0,\delta)$ in
$\mathbb{R}^3$, and put
\[
U_\delta=\{y\in\mathbb{R}^3:\delta<\|y\|<1\},
\qquad F_\delta=\chi_{U_\delta}.
\]

\begin{proposition}
\label{three_dimensional_shell_obstruction}
For every $0<\delta<1$, there is an $x_\delta>1$ such that
$r=x_\delta$ is the unique maximizing radius for
$A_{F_\delta}(x_\delta,r)$. One can take $x_\delta$ to be the larger
root of
\begin{equation}
\label{three_dimensional_shell_root}
4(1+\delta)x^2-8(1+\delta+\delta^2)x
       +3(1+\delta)(1+\delta^2)=0.
\end{equation}
In particular,
\[
x_{1/2}=\frac{14+\sqrt{61}}{12}.
\]
Moreover, for any $\epsilon>0$ and $0<\sigma<a$ with $a+\sigma<\delta$,
the two functions
\[
f_+=F_\delta+\epsilon\chi_{B(ae_1,\sigma)},
\qquad
f_-=F_\delta+\epsilon\chi_{B(-ae_1,\sigma)}
\]
have different maximal functions at $X=x_\delta e_1$.
\end{proposition}

\begin{proof}
For $x>1$, put $A(x,r)=A_{F_\delta}(x,r)$. The three-dimensional
intersection formula \eqref{vol_3_balls_inter} gives the following
derivatives in the three partial-overlap regimes:
\begin{equation}
\label{three_dimensional_shell_derivatives}
\frac{\partial A}{\partial r}(x,r)=
\begin{cases}
\displaystyle
\frac{3\bigl(r^2-(x-1)^2\bigr)
              \bigl(x^2+2x-3-r^2\bigr)}{16xr^4},
 &x-1<r<x-\delta,\\[8pt]
\displaystyle
\frac{3(1-\delta)}{16xr^4}
       \bigl(C_\delta(x)-2(1+\delta)r^2\bigr),
 &x-\delta<r<x+\delta,\\[8pt]
\displaystyle -\frac{3}{16xr^4}K_\delta(x,r),
 &x+\delta<r<x+1,
\end{cases}
\end{equation}
where
\begin{align*}
C_\delta(x)
 &=6(1+\delta)x^2-8(1+\delta+\delta^2)x
       +3(1+\delta)(1+\delta^2),\\
K_\delta(x,r)
 &=r^4-2x^2r^2+2r^2+x^4-6x^2+8x-3-16x\delta^3.
\end{align*}
The expressions for $\partial A/\partial r$ agree at their common
endpoints. Also $A(x,r)=0$ for $r\leq x-1$, whereas
$A(x,r)=(1-\delta^3)/r^3$ for $r\geq x+1$.

The polynomial on the left of \eqref{three_dimensional_shell_root}
has positive leading coefficient and takes the value
\[
3\delta^3-5\delta^2-\delta-1<0
\]
at $x=1$. It therefore has a unique root $x_\delta>1$.
At this value of $x$, we have
$C_\delta(x)=2(1+\delta)x^2$. Hence the middle expression in
\eqref{three_dimensional_shell_derivatives} is positive for $r<x$
and negative for $r>x$.

In the first regime, the sign of the derivative is the sign of the
strictly decreasing factor $x^2+2x-3-r^2$. Since the derivative is
positive at $r=x-\delta$, it is positive throughout that regime.
In the last regime,
\[
\frac{\partial K_\delta}{\partial r}(x,r)
   =4r(r^2-x^2+1)>0.
\]
The derivative of $A$ is negative at $r=x+\delta$, so it remains
negative up to $r=x+1$ and is negative thereafter. Thus $A(x,\cdot)$
is strictly increasing on $(x-1,x)$ and strictly decreasing on
$(x,\infty)$, proving the assertion about the maximizing radius.

Set $X=x_\delta e_1$ and $m=M_cF_\delta(X)$. The ball
$B(X,x_\delta)$ contains $B(ae_1,\sigma)$ entirely. Therefore
\[
M_cf_+(X)\geq m+\epsilon\frac{\sigma^3}{x_\delta^3}.
\]
In contrast, an averaging ball centered at $X$ does not meet
$B(-ae_1,\sigma)$ until its radius exceeds $x_\delta+a-\sigma$.
Its contribution from this added ball is consequently at most
$\epsilon\sigma^3/(x_\delta+a-\sigma)^3$. Since its contribution
from the shell is at most $m$, we obtain
\begin{equation}
\label{three_dimensional_moving_ball_obstruction}
\begin{aligned}
M_cf_-(X)
&\leq m+\epsilon\frac{\sigma^3}{(x_\delta+a-\sigma)^3}\\
&<m+\epsilon\frac{\sigma^3}{x_\delta^3}
\leq M_cf_+(X).
\end{aligned}
\end{equation}
This holds for every $\epsilon>0$, however small.
\end{proof}

Both added balls in the proposition can be placed in an arbitrarily
small neighborhood of the origin. Thus a single uniform shell does
not allow the translation-invariant family used in the planar
construction. This conclusion concerns that particular construction;
it does not rule out other backgrounds or other pairs of functions.

\subsection{Radial averages and the choice of three layers}

The underlying principle is unchanged. We seek a background for which,
after the perturbations have been added, the maximal function can be
computed using only balls that either miss a small central ball or
contain it entirely. Equal total masses inside that central ball then
produce the same averages at all relevant radii. The additional task
is to check that the perturbations do not create a new maximizing
radius among the excluded balls.

Radial functions make this verification explicit in dimension three.
Let $W(y)=w(\|y\|)$ be locally bounded and radial, and put
\[
I_W(x,r)=\int_{B(xe_1,r)}W(y)\,dy,
\qquad p=x+r,\qquad q=|x-r|,\qquad x>0.
\]
The portion of the sphere $\|y\|=s$ inside $B(xe_1,r)$ has area
\[
\frac{\pi s}{x}\bigl(r^2-(x-s)^2\bigr)
\]
when $q<s<p$. Spheres with $s<(r-x)_+$ are wholly contained in
the averaging ball. Integrating these spherical sections therefore gives
\begin{equation}
\label{three_dimensional_radial_integral}
\begin{aligned}
I_W(x,r)
={}&4\pi\int_0^{(r-x)_+}s^2w(s)\,ds\\
&+\frac{\pi}{x}\int_q^p
        s\bigl(r^2-(x-s)^2\bigr)w(s)\,ds,
\end{aligned}
\end{equation}
where $(r-x)_+=\max\{r-x,0\}$. Integrating on the boundary of the
averaging ball, or differentiating this identity, also gives
\begin{equation}
\label{three_dimensional_radial_derivative}
\frac{\partial I_W}{\partial r}(x,r)
     =\frac{2\pi r}{x}\int_q^p s w(s)\,ds.
\end{equation}
The right-hand side is continuous in $r>0$ even when $w$ has jumps.
Thus $I_W(x,\cdot)$ and $A_W(x,\cdot)$ are continuously differentiable.
If
\[
D_W(x,r)=r\frac{\partial I_W}{\partial r}(x,r)-3I_W(x,r),
\]
then
\begin{equation}
\label{three_dimensional_average_derivative}
\frac{\partial A_W}{\partial r}(x,r)
       =\frac{3}{4\pi r^4}D_W(x,r).
\end{equation}
In particular, a maximum at a positive radius requires $D_W(x,r)=0$.
For the piecewise smooth profiles below, $D_W(x,\cdot)$ is locally
absolutely continuous, so its derivative inequalities can also be
integrated across the finitely many breakpoints.

We choose the background in the form
\[
F(y)=\frac{u(\|y\|)}{\|y\|}.
\]
The factor $1/\|y\|$ cancels the factor $s$ in the second integral
in \eqref{three_dimensional_radial_integral}. The resulting kernel
\[
r^2-(x-s)^2=(r^2-x^2)+2xs-s^2
\]
is quadratic in $s$. This is an exact integration formula, not an
application of the mean-value property for harmonic functions; nor
does it require an explicit formula for the supremum over $r$.

For example, suppose that $u$ vanishes on $[0,q]$. Set
\[
U_j(p)=\int_0^p s^j u(s)\,ds,\qquad j=0,1,2.
\]
Then \eqref{three_dimensional_radial_integral} and
\eqref{three_dimensional_radial_derivative} become
\begin{align}
I_F(x,r)
 &=\frac{\pi}{x}\bigl((r^2-x^2)U_0(p)
                    +2xU_1(p)-U_2(p)\bigr),
       \label{three_dimensional_moment_integral}\\
\frac{x}{\pi}D_F(x,r)
 &=(3x^2-r^2)U_0(p)-6xU_1(p)+3U_2(p).
       \label{three_dimensional_moment_derivative}
\end{align}
The full-sphere contribution is zero here because $(r-x)_+\leq q$.

The three layers used below are centered at radii $1$, $5/3$, and $6$,
with respective weights $1$, $3$, and $1$ in the profile $u$.
To explain the choice, first consider the zero-width limit and set
$r=x$. As $2x$ passes the first, second, and third layers, the
right-hand side of \eqref{three_dimensional_moment_derivative}
becomes, respectively,
\[
2x^2-6x+3,\qquad
8x^2-36x+28,\qquad
10x^2-72x+136.
\]
The last polynomial is strictly positive. In the ranges where the
first two polynomials apply, their only zeros occur at
$x=(3-\sqrt3)/2$ and $x=1$. Near the first of these centers, a larger
averaging ball reaching the second layer has a higher average. Near
the second, a small averaging ball meeting the first layer has a
higher average. The proof below makes these comparisons uniform for
layers of positive thickness and for the two perturbed functions.

Although the quadratic kernel suggests cancelling three moments of
a radial perturbation, that is not needed here. We exclude every
averaging sphere that crosses the small central ball. The cancellation
condition on the perturbation is therefore just its zero total integral.

\subsection{An explicit positive pair}

Fix $\eta=10^{-3}$ and define, for $s\geq0$,
\begin{equation}
\label{three_dimensional_profile}
u(s)=\frac{1}{2\eta}\Bigl(
 \chi_{(1-\eta,1+\eta)}(s)
 +3\chi_{(5/3-\eta,5/3+\eta)}(s)
 +\chi_{(6-\eta,6+\eta)}(s)\Bigr).
\end{equation}
Put $F(0)=0$, $F(y)=u(\|y\|)/\|y\|$ for $y\ne0$, and
\begin{equation}
\label{three_dimensional_perturbation}
h=\frac12\chi_{B_{\eta/2}}
          -\frac1{14}\chi_{B_\eta\setminus B_{\eta/2}},
\qquad B_t=B(0,t).
\end{equation}
\begin{samepage}
\begin{theorem}
\label{three_dimensional_theorem}
Let $F$ and $h$ be defined above. For every $0<\epsilon\leq10^{-6}$,
the functions
\begin{equation}
\label{three_dimensional_pair}
\begin{aligned}
f(y)&=F(y)+\epsilon\bigl(e^{-\|y\|}+h(y)\bigr),\\
g(y)&=F(y)+\epsilon\bigl(e^{-\|y\|}-h(y)\bigr)
\end{aligned}
\end{equation}
are bounded, radial, strictly positive functions in $L^1(\mathbb{R}^3)$
such that
\begin{equation}
\label{three_dimensional_conclusion}
M_cf(x)=M_cg(x)\quad\text{for every }x\in\mathbb{R}^3,
\qquad \|f-g\|_1=\frac{\epsilon\pi\eta^3}{3}>0.
\end{equation}
Their common $L^1$ norm is $48\pi+8\pi\epsilon$.
\end{theorem}
\end{samepage}

\begin{proof}
The volume of $B_\eta\setminus B_{\eta/2}$ is seven times that of
$B_{\eta/2}$. Hence
\begin{equation}
\label{three_dimensional_zero_mass}
\int_{\mathbb{R}^3}h=0,
\qquad \|h\|_1=\frac{\pi\eta^3}{6}.
\end{equation}
On $B_\eta$, we have $e^{-\|y\|}\geq e^{-\eta}>1/2\geq|h(y)|$;
off $B_\eta$, the exponential is positive and $h=0$. Thus both
functions in \eqref{three_dimensional_pair} are strictly positive.
Also
\[
\int_0^\infty s u(s)\,ds=1+3\cdot\frac53+6=12,
\qquad
\int_{\mathbb{R}^3}e^{-\|y\|}\,dy=8\pi.
\]
This proves integrability and the asserted norm formulas. Boundedness
follows since all three layers are separated from the origin.

\medskip
\noindent\textit{The radii to be excluded.}
We will prove for each $\phi\in\{f,g\}$ that
\begin{equation}
\label{three_dimensional_safe_radii}
M_c\phi(xe_1)=\sup_{\substack{r>0\\ |r-x|>\eta}}A_\phi(x,r),
\qquad x\geq0.
\end{equation}
Call a positive radius excluded if $|r-x|\leq\eta$.
At every other radius, $B(xe_1,r)$ either misses $B_\eta$ or contains
it entirely. The averages of $f$ and $g$ are equal at all these
other radii by \eqref{three_dimensional_zero_mass}.
It therefore suffices to establish \eqref{three_dimensional_safe_radii}.

For the derivative estimates assume first that $x>0$. Let
\[
p=x+r,\qquad q=|x-r|,\qquad t=r-x,
\]
and use $U_j(p)$ from \eqref{three_dimensional_moment_integral}.
For an excluded radius, $q\leq\eta<1-\eta$, so both moment formulas
apply. Write
\[
H(x,r)=\frac{x}{\pi}D_F(x,r)
       =(3x^2-r^2)U_0(p)-6xU_1(p)+3U_2(p).
\]
For reference, the moments between the layers are
\begin{equation}
\label{three_dimensional_moments}
\bigl(U_0(p),U_1(p),3U_2(p)\bigr)=
\begin{cases}
(1,1,3+\eta^2),&1+\eta<p<5/3-\eta,\\
(4,6,28+4\eta^2),&5/3+\eta<p<6-\eta,\\
(5,12,136+5\eta^2),&p>6+\eta.
\end{cases}
\end{equation}
These identities follow by integration over the three symmetric
intervals in \eqref{three_dimensional_profile}.

Let $V_\pm(y)=e^{-\|y\|}\pm h(y)$ and
$V_\pm(y)=v_\pm(\|y\|)$. Then
\[
0<v_\pm\leq\frac32,\qquad
\int_0^\infty s v_\pm(s)\,ds
          =1\pm\frac{\eta^2}{28}<2,
\qquad \int_{\mathbb{R}^3}V_\pm=8\pi.
\]
By \eqref{three_dimensional_radial_derivative},
\[
\left|\frac{x}{\pi}D_{V_\pm}(x,r)\right|
   \leq4r^2+24x\leq20(1+x^2)
       \qquad (|r-x|\leq\eta).
\]
For the last inequality, use $r\leq x+\eta$,
$4r^2\leq8x^2+8\eta^2$, and $24x\leq12(1+x^2)$.
Consequently,
\begin{equation}
\label{three_dimensional_derivative_error}
\frac{x}{\pi}D_\phi(x,r)=H(x,r)+\epsilon E_\pm(x,r),
\qquad |E_\pm(x,r)|\leq20(1+x^2).
\end{equation}
We now show that no excluded radius can be a global maximizer.

\medskip
\noindent\textit{Case 1: $p\leq1-\eta$.}
Here the averaging ball misses all three layers and $x\leq1/2$.
Thus $A_\phi(x,r)\leq3\epsilon/2$. The radius
$R=x+6+\eta$ contains the entire support of $F$ and gives
\[
A_\phi(x,R)\geq A_F(x,R)
       =\frac{36}{(x+6+\eta)^3}>\frac{13}{100}>\frac{3\epsilon}{2}.
\]
This also treats $x=0$.

\medskip
\noindent\textit{Case 2: $p$ lies in a layer or at one of its endpoints.}
Inside a layer, $u(p)\geq500$ and $U_0(p)\leq5$.
Differentiating \eqref{three_dimensional_radial_derivative} gives
\[
\frac{\partial D_F}{\partial r}(x,r)
       =\frac{2\pi r}{x}\bigl(r u(p)-U_0(p)\bigr).
\]
For $V_\pm$, the corresponding identity, almost everywhere, is
\begin{align*}
\frac{\partial D_{V_\pm}}{\partial r}(x,r)
 =\frac{2\pi r}{x}\Bigg[&
 r\bigl(pv_\pm(p)-\operatorname{sgn}(r-x)qv_\pm(q)\bigr)\\
 &-\int_q^p s v_\pm(s)\,ds\Bigg].
\end{align*}
To control a full neighborhood of an excluded radius, we may allow
$q\leq2\eta$. If $p$ is still in a layer, then
\[
\frac{249}{500}<r<\frac{1501}{500},\qquad
qv_\pm(q)\leq3\eta.
\]
The expression in square brackets is therefore greater than
$-2-3\eta r>-201/100$. It follows that
\begin{equation}
\label{three_dimensional_layer_convexity}
\frac{\partial D_\phi}{\partial r}(x,r)
 >\frac{2\pi r}{x}\left(244-\frac{201}{100}\epsilon\right)>0.
\end{equation}
Thus a zero of $D_\phi$ in the interior of a layer gives a strict
local minimum of $A_\phi$, not a maximum. If $p$ is at the entrance
to a layer and $D_\phi=0$, the average increases immediately to the
right. At the exit, a zero of $D_\phi$ is approached through negative
values from the left, so the preceding averages are larger.
The one-sided assertions follow by integrating
\eqref{three_dimensional_layer_convexity} within the layer.
They remain valid at jumps of $v_\pm$. If $D_\phi\ne0$, continuous
differentiability already excludes a local maximum. This treats
all layer endpoints as well as their interiors.

\medskip
\noindent\textit{Case 3: $1+\eta<p<5/3-\eta$.}
Here $1/2<x<5/6$, and \eqref{three_dimensional_moments} gives
\[
H(x,r)=2x^2-6x+3-2xt-t^2+\eta^2.
\]
Since $|t|\leq\eta$, the difference between
$xD_\phi(x,r)/\pi$ and $P_1(x)=2x^2-6x+3$ is at most
\[
\frac53\eta+\eta^2
       +20\epsilon\left(1+\frac{25}{36}\right)<\frac1{500}.
\]
The polynomial $P_1$ is strictly decreasing on $(1/2,5/6)$ and
\[
P_1(63/100)=\frac{69}{5000},\qquad
P_1(16/25)=-\frac{13}{625}.
\]
Thus $D_\phi(x,r)=0$ is possible only when
\begin{equation}
\label{three_dimensional_first_critical_range}
\frac{63}{100}<x<\frac{16}{25}.
\end{equation}
In this range, \eqref{three_dimensional_moment_integral} implies
\begin{align*}
A_F(x,r)
 &=\frac{3}{4xr^3}
       \left(2xt+t^2+2x-1-\frac{\eta^2}{3}\right)\\
 &\leq\frac{3\bigl(2(16/25)\eta+\eta^2+7/25\bigr)}
             {4(63/100)(63/100-\eta)^3}
 <\frac{27}{20}.
\end{align*}
On the other hand, take $R=3/2$. This ball meets both of the first
two layers, with $|R-x|<1-\eta$ and $x+R>5/3+\eta$.
Consequently,
\[
A_F(x,R)=\frac{3}{4xR^3}
       \left(4R^2-4x^2+12x-\frac{28}{3}-\frac{4\eta^2}{3}\right).
\]
The expression in parentheses is increasing throughout
\eqref{three_dimensional_first_critical_range}. Replacing $x$ there
by $63/100$ and the $x$ in the denominator by $16/25$ gives
\[
A_F(x,R)\geq
\frac{3\bigl(4R^2-4(63/100)^2+12(63/100)-28/3-4\eta^2/3\bigr)}
     {4(16/25)R^3}
>\frac{39}{20}.
\]
Since $0<V_\pm\leq3/2$, every possible stationary excluded radius
in this case satisfies
\[
A_\phi(x,r)<\frac{27}{20}+\frac{3\epsilon}{2}
       <\frac{39}{20}<A_\phi(x,R).
\]
It cannot be a global maximizer.

\medskip
\noindent\textit{Case 4: $5/3+\eta<p<6-\eta$.}
Now $5/6<x<3$ and
\[
H(x,r)=8x^2-36x+28-8xt-4t^2+4\eta^2.
\]
The difference between $xD_\phi(x,r)/\pi$ and
\[
P_2(x)=8x^2-36x+28=4(x-1)(2x-7)
\]
is smaller than
\[
24\eta+4\eta^2+200\epsilon<\frac1{40}.
\]
On $(5/6,99/100]$, $P_2\geq P_2(99/100)=251/1250$.
On $[101/100,3)$, the maximum of this convex quadratic is at one
of the endpoints and is at most
$P_2(101/100)=-249/1250$. Hence stationarity requires
\begin{equation}
\label{three_dimensional_second_critical_range}
\frac{99}{100}<x<\frac{101}{100}.
\end{equation}
In this range,
\begin{align*}
A_F(x,r)
 &=\frac{3}{4xr^3}
       \left(8xt+4t^2+12x-\frac{28}{3}-\frac{4\eta^2}{3}\right)\\
 &\leq\frac{3\bigl(8(101/100)\eta+4\eta^2+12(101/100)-28/3\bigr)}
              {4(99/100)(99/100-\eta)^3}
 <\frac{11}{5}.
\end{align*}
Take instead $R=1/10$. This ball meets the first layer only, and
its radial range contains the whole interval $(1-\eta,1+\eta)$.
Formula \eqref{three_dimensional_radial_integral} yields
\begin{align*}
A_F(x,R)
 &=\frac{3}{4xR^3}
       \left(R^2-(x-1)^2-\frac{\eta^2}{3}\right)\\
 &\geq\frac{3\bigl(R^2-10^{-4}-\eta^2/3\bigr)}
              {4(101/100)R^3}
 >\frac{73}{10}.
\end{align*}
Thus again
\[
A_\phi(x,r)<\frac{11}{5}+\frac{3\epsilon}{2}
       <\frac{73}{10}<A_\phi(x,R),
\]
which excludes every possible stationary radius in this case.

\medskip
\noindent\textit{Case 5: $p>6+\eta$.}
All three layers contribute, and
\begin{align*}
H(x,r)
 &=10x^2-(72+10t)x+136+5\eta^2-5t^2\\
 &\geq10x^2-\frac{7201}{100}x+136
 >\frac{1+x^2}{10}.
\end{align*}
For the last inequality, subtract $(1+x^2)/10$; the resulting
quadratic has positive leading coefficient and discriminant
\[
\left(\frac{7201}{100}\right)^2
       -4\cdot\frac{99}{10}\cdot\frac{1359}{10}
       =-\frac{1961999}{10000}<0.
\]
Together with \eqref{three_dimensional_derivative_error}, this gives
\[
\frac{x}{\pi}D_\phi(x,r)
       >\left(\frac1{10}-20\epsilon\right)(1+x^2)>0.
\]
The average is increasing at this radius, so no maximum is possible.

\medskip
\noindent\textit{Taking the supremum and completing the proof.}
The five cases show that no positive excluded radius is a global
maximizer. If $x>\eta$, the excluded radii form the compact interval
$[x-\eta,x+\eta]\subset(0,\infty)$, on which $A_\phi(x,\cdot)$ is
continuous. If its maximum on this interval were at least the
supremum over the other radii, an excluded radius would be a global
maximizer. This has been ruled out, proving
\eqref{three_dimensional_safe_radii} for $x>\eta$.

If $0\leq x\leq\eta$, every excluded ball lies in $B(0,3\eta)$,
where $F=0$, so its average is at most $3\epsilon/2$. Case~1 supplies
a larger average at the nonexcluded radius $x+6+\eta$. This also
proves \eqref{three_dimensional_safe_radii} when a supremum over
excluded radii could be approached as $r\downarrow0$.

Finally, for $|r-x|>\eta$, the averaging ball either misses $B_\eta$
or contains it completely. Equation \eqref{three_dimensional_zero_mass}
therefore gives $A_f(x,r)=A_g(x,r)$. Taking the suprema in
\eqref{three_dimensional_safe_radii} proves the equality of the maximal
functions on the ray $xe_1$. Since $f$ and $g$ are radial, it holds
at every point of $\mathbb{R}^3$.
\end{proof}

\begin{remark}
The construction changes the radial distribution of mass inside
$B_\eta$, rather than translating a single bump. The common
exponential term makes both functions strictly positive on all of
$\mathbb{R}^3$. Its effect on the maximizing radii is included in the
proof; adding an arbitrary common positive function would not, by
itself, preserve equality of maximal functions.
\end{remark}

\subsection{The maximal function of the unit ball}

The same intersection formula also gives a short explicit expression
for the central maximal function of a three-dimensional ball.

\begin{corollary}
For $U=B(0,1)\subset\mathbb{R}^3$,
\begin{equation}
\label{maximal_ball}
M_c(\chi_U)(x)=
\begin{cases}
1,&\|x\|<1,\\[2pt]
\displaystyle\frac12,&\|x\|=1,\\[6pt]
\displaystyle\frac12-\frac12
 \sqrt{\frac{\|x\|-1}{\|x\|+3}}
 \left(1+\frac{2}{\|x\|}\right),&\|x\|>1.
\end{cases}
\end{equation}
For $\|x\|>1$ the unique maximizing radius is
\[
r_x=\sqrt{(\|x\|-1)(\|x\|+3)}.
\]
\end{corollary}

\begin{proof}
For $x>1$ and $x-1<r<x+1$, formula \eqref{vol_3_balls_inter} gives
\begin{equation}
\label{average}
A_{\chi_U}(x,r)
 =\frac{(r-x+1)^2
       \bigl(x^2+2x(r+1)-3(r-1)^2\bigr)}{16xr^3}.
\end{equation}
Its derivative is the first expression in
\eqref{three_dimensional_shell_derivatives}, throughout this whole
partial-overlap interval. It changes sign from positive to negative
at $r=\sqrt{x^2+2x-3}$, which lies strictly between $x-1$ and $x+1$.
Before that interval the average is zero; after it the average is
$r^{-3}$ and is decreasing. Substitution of the maximizing radius
into \eqref{average} proves \eqref{maximal_ball} for exterior centers.

For $x=1$ and $0<r<2$, \eqref{average} reduces to
\[
A_{\chi_U}(1,r)=\frac12-\frac{3r}{16}.
\]
For $r\geq2$ the average is $r^{-3}\leq1/8$. Thus the supremum
is $1/2$, approached as $r\downarrow0$. For centers inside $U$,
sufficiently small balls have average $1$, which is also the upper
bound for every average.
\end{proof}

In particular, the exterior boundary asymptotic is
\begin{equation}
\label{three_dimensional_boundary_asymptotic}
\frac12-M_c(\chi_U)(x)\sim\frac34(\|x\|-1)^{1/2}
       \qquad\text{as }\|x\|\downarrow1.
\end{equation}

\begin{remark}
This boundary asymptotic is related to the level-set geometry underlying
\emph{Solyanik estimates}; see Hagelstein and
Parissis~\cite{hagelsteinparissis2015}. To make the connection precise
for the fixed ball $U$, put
\[
H_U(\alpha)=\frac{|\{x\in\mathbb{R}^3:
                         M_c(\chi_U)(x)>\alpha\}|}{|U|}.
\]
For $0<\alpha<1/2$, the level set is the ball of radius $R_\alpha>1$,
where \eqref{maximal_ball} gives
\[
R_\alpha^2(R_\alpha+3)=\frac{1}{\alpha(1-\alpha)}.
\]
Consequently,
\begin{equation}
\label{three_dimensional_ball_halo}
H_U(\alpha)-1\sim\frac{16}{3}\left(\frac12-\alpha\right)^2
                   \qquad\text{as }\alpha\uparrow\frac12,
\end{equation}
whereas $H_U(\alpha)=1$ for $1/2\leq\alpha<1$.
Indeed, writing $\alpha=1/2-t$ in the preceding equation gives
$R_\alpha-1\sim16t^2/9$, and $H_U(\alpha)=R_\alpha^3$.
Classical Solyanik estimates concern sharp Tauberian constants, obtained
by taking a supremum over measurable sets of finite positive measure,
as the level tends to $1$. Formula \eqref{three_dimensional_ball_halo}
is a fixed-set boundary analogue at level $1/2$, rather than such a
uniform estimate at level $1$.
\end{remark}


\section{The $n$-Dimensional Result for $n\geq4$}
\label{n_dimensional_result}

We now give a construction in every dimension $n\geq4$. The numerical
parameters of the preceding three-layer example are replaced by
parameters depending on $n$. The new background consists of only two
uniform concentric shells.

The idea is to start with $\chi_{U_\delta}$ and locate the unique
exterior center at which the radius $r=x$ is stationary for its
average. We then add a thin outer shell and choose its position so
that the corresponding stationary center lies in the interior of
this new shell. At that center, sufficiently small balls have average
$1$, whereas the ball of radius $r=x$ meets the central hole and has
average strictly less than $1$. At every other exterior center, the
radius $r=x$ is not stationary. The position of the added shell must
be adjusted to account for the change that the shell itself produces
in the stationary center.

After this adjustment, a uniform interval of radii around $r=x$ can
be excluded, even after sufficiently small perturbations. As before,
a zero-integral redistribution of mass inside a smaller central ball
then leaves the maximal function unchanged.

\begin{theorem}
\label{n_dimensional_theorem}
For every integer $n\geq4$, put $\delta=1/(4n)$. There exist
$\tau>0$, $\lambda>1+\tau$, $0<\eta<\delta$, and $\epsilon_0>0$
with the following property. Set
\begin{equation}
\label{n_dimensional_background}
F(y)=\chi_{U_\delta}(y)
 +\chi_{\{\lambda-\tau<\|y\|<\lambda+\tau\}}(y),
\qquad U_\delta=\{\delta<\|y\|<1\},
\end{equation}
and
\begin{equation}
\label{n_dimensional_perturbation}
h_n=\frac12\chi_{B_{\eta/2}}
 -\frac{1}{2(2^n-1)}\chi_{B_\eta\setminus B_{\eta/2}},
\qquad B_s=B(0,s).
\end{equation}
For every $0<\epsilon\leq\epsilon_0$, the functions
\begin{equation}
\label{n_dimensional_pair}
\begin{aligned}
f(y)&=F(y)+\epsilon\bigl(e^{-\|y\|}+h_n(y)\bigr),\\
g(y)&=F(y)+\epsilon\bigl(e^{-\|y\|}-h_n(y)\bigr)
\end{aligned}
\end{equation}
are bounded, radial, strictly positive, and integrable on
$\mathbb{R}^n$. Moreover,
\begin{equation}
\label{n_dimensional_conclusion}
M_cf(X)=M_cg(X)\quad(X\in\mathbb{R}^n),
\qquad
\|f-g\|_1=2\epsilon c_n\left(\frac{\eta}{2}\right)^n>0.
\end{equation}
\end{theorem}

We prepare the proof by deriving the diagonal derivative of a ball
average, choosing the second shell, and establishing stability of
the excluded interval of radii.

\subsection{Diagonal derivatives of radial averages}

For a bounded radial function $W$, write
\[
I_W(x,r)=\int_{B(xe_1,r)}W(y)\,dy,
\qquad A_W(x,r)=\frac{I_W(x,r)}{c_nr^n},\qquad x>0.
\]
Whenever the derivative exists, put
\[
\mathcal D_W(x)=x\frac{\partial I_W}{\partial r}(x,x)-nI_W(x,x).
\]
Thus
\begin{equation}
\label{n_dimensional_diagonal_derivative}
\frac{\partial A_W}{\partial r}(x,x)
       =\frac{\mathcal D_W(x)}{c_nx^{n+1}}.
\end{equation}
In particular, a nonzero value of $\mathcal D_W(x)$ excludes $r=x$
as a maximizing radius.

Define
\[
Q_n(a)=\int_a^1(1-t^2)^{(n-3)/2}\,dt\qquad(0\leq a\leq1)
\]
and
\begin{equation}
\label{n_dimensional_phi}
\Phi_n(t)=
\begin{cases}
0,&0\leq t\leq1/2,\\[4pt]
\displaystyle
\frac{t}{n-1}\left(1-\frac{1}{4t^2}\right)^{(n-1)/2}
 -Q_n\!\left(\frac{1}{2t}\right),&t>1/2.
\end{cases}
\end{equation}
The function $\Phi_n$ is continuously differentiable on $[0,\infty)$
for $n\geq4$. Direct differentiation gives
\begin{equation}
\label{n_dimensional_phi_derivative}
\Phi_n'(t)=\frac{1}{n-1}
 \left(1-\frac{1}{4t^2}\right)^{(n-3)/2}
 \left(1-\frac{n}{4t^2}\right),\qquad t>1/2.
\end{equation}
Its derivative tends to zero as $t\downarrow1/2$, agreeing with the
zero derivative on the other side.

\begin{lemma}
\label{n_dimensional_ball_derivative}
For $s>0$ and $x>0$,
\begin{equation}
\label{n_dimensional_ball_kernel}
\mathcal D_{\chi_{B_s}}(x)
       =(n-1)c_{n-1}s^n\Phi_n(x/s).
\end{equation}
\end{lemma}

\begin{proof}
If $2x<s$, the averaging ball remains inside $B_s$ for radii near
$r=x$, so its normalized average is $1$ and the diagonal derivative
is zero. Suppose that $2x>s$, and put $\kappa_n=(n-1)c_{n-1}$,
the surface area of the unit sphere in $\mathbb{R}^{n-1}$.
At distance $v$ from the origin, the angular condition for membership
in $B(xe_1,x)$ is $\cos\theta>v/(2x)$. Polar integration and
differentiation with respect to the averaging radius therefore give
\begin{align*}
I_{\chi_{B_s}}(x,x)
 &=\kappa_n\int_0^s v^{n-1}Q_n\!\left(\frac{v}{2x}\right)\,dv,\\
\frac{\partial I_{\chi_{B_s}}}{\partial r}(x,x)
 &=\kappa_n\int_0^s v^{n-2}
          \left(1-\frac{v^2}{4x^2}\right)^{(n-3)/2}\,dv.
\end{align*}
In the second identity one may differentiate the angular threshold
$(x^2+v^2-r^2)/(2xv)$ before setting $r=x$.
Integrating the first identity by parts yields
\begin{align*}
\frac{\mathcal D_{\chi_{B_s}}(x)}{\kappa_n}
={}&x\int_0^s v^{n-2}
 \left(1-\frac{v^2}{2x^2}\right)
 \left(1-\frac{v^2}{4x^2}\right)^{(n-3)/2}\,dv\\
&-s^nQ_n\!\left(\frac{s}{2x}\right).
\end{align*}
The integrand in the first line, including the factor $x$, is the
derivative with respect to $v$ of
\[
\frac{x}{n-1}v^{n-1}
       \left(1-\frac{v^2}{4x^2}\right)^{(n-1)/2}.
\]
Evaluation at $v=s$ proves \eqref{n_dimensional_ball_kernel}.
At $2x=s$, the first radial derivative of the intersection volume
is continuous at internal tangency. Equivalently, the volume lost
at tangency is of higher than first order in the change of radius.
Thus the same identity holds there, with both sides equal to zero.
\end{proof}

In particular, for $F_0=\chi_{U_\delta}$,
\begin{equation}
\label{n_dimensional_psi}
\mathcal D_{F_0}(x)=(n-1)c_{n-1}\Psi_n(x),
\qquad
\Psi_n(x)=\Phi_n(x)-\delta^n\Phi_n(x/\delta).
\end{equation}

\begin{lemma}
\label{n_dimensional_unique_root}
For $\delta=1/(4n)$, the function $\Psi_n$ has exactly one zero
$x_n$ on $[1,\infty)$. This zero satisfies $x_n>1$ and
$\Psi_n'(x_n)>0$.
\end{lemma}

\begin{proof}
First,
\[
\Phi_n(1)=-\int_{1/2}^1(1-t)(1-t^2)^{(n-3)/2}\,dt<0.
\]
Also $Q_n(1/(8n))\leq1$, and Bernoulli's inequality gives
\[
\Phi_n(4n)
 \geq\frac{4n}{n-1}\left(1-\frac{1}{64n}\right)-1>0.
\]
Indeed, $(n-1)/2\leq n$ and
$(1-1/(64n^2))^n\geq1-1/(64n)$.
It follows that $\Psi_n(1)<0$.
As $x\to\infty$, \eqref{n_dimensional_phi} gives
\begin{equation}
\label{n_dimensional_psi_infinity}
\Psi_n(x)=\frac{1-\delta^{n-1}}{n-1}x
 -(1-\delta^n)Q_n(0)+O(x^{-1}).
\end{equation}
In particular, $\Psi_n(x)\to\infty$.

For $x\geq1$, we have $x/\delta\geq4n>\sqrt n/2$, so
$\Phi_n'(x/\delta)>0$. On $1\leq x\leq\sqrt n/2$,
\eqref{n_dimensional_phi_derivative} implies $\Psi_n'(x)<0$.
For $x>\sqrt n/2$, put $z=1/(4x^2)$. Then
\begin{equation}
\label{n_dimensional_derivative_ratio}
\frac{\Phi_n'(x)}{\Phi_n'(x/\delta)}
 =\left(\frac{1-z}{1-\delta^2z}\right)^{(n-3)/2}
       \frac{1-nz}{1-n\delta^2z}.
\end{equation}
Each factor on the right is positive and strictly decreases with
$z\in(0,1/n)$. Since $z$ decreases with $x$, the ratio strictly
increases from $0$ to $1$ as $x$ increases from $\sqrt n/2$ to
infinity. Hence
\[
\Psi_n'(x)=\Phi_n'(x/\delta)
 \left(\frac{\Phi_n'(x)}{\Phi_n'(x/\delta)}-\delta^{n-1}\right)
\]
changes sign exactly once, from negative to positive.
Thus $\Psi_n$ first strictly decreases and then strictly increases.
Its negative value at $1$ and \eqref{n_dimensional_psi_infinity}
prove the assertion, including simplicity of the zero.
\end{proof}

\subsection{Positioning the second shell}

\begin{lemma}
\label{n_dimensional_matched_shell}
There exist $\tau>0$ and $\lambda>1+\tau$ such that the background
$F$ in \eqref{n_dimensional_background} satisfies
\begin{equation}
\label{n_dimensional_strict_gap}
A_F(x,x)<M_cF(xe_1)\qquad\text{for every }x>0.
\end{equation}
The parameters can be chosen with $\tau$ arbitrarily small and
$\lambda\to x_n$ as $\tau\downarrow0$.
\end{lemma}

\begin{proof}
Let $K(s,x)=s^n\Phi_n(x/s)$, and set
\[
R_{\tau,\lambda}(x)=K(\lambda+\tau,x)-K(\lambda-\tau,x).
\]
By \eqref{n_dimensional_ball_kernel},
\begin{equation}
\label{n_dimensional_shell_derivative}
\mathcal D_F(x)=(n-1)c_{n-1}
                \bigl(\Psi_n(x)+R_{\tau,\lambda}(x)\bigr).
\end{equation}
We choose $\lambda$ by the matching equation
\begin{equation}
\label{n_dimensional_matching_equation}
\Psi_n(\lambda)+R_{\tau,\lambda}(\lambda)=0.
\end{equation}
At $\tau=0$ this is $\Psi_n(\lambda)=0$.
The left-hand side is smooth near $(\tau,\lambda)=(0,x_n)$;
its derivative with respect to $\lambda$ at that point is
$\Psi_n'(x_n)>0$. The implicit function theorem therefore gives
a function $\lambda=\lambda(\tau)$ for all sufficiently small
$|\tau|$, with $\lambda(0)=x_n$, satisfying
\eqref{n_dimensional_matching_equation}. For small $\tau>0$,
we have $\lambda(\tau)-\tau>1$.

We verify that no other exterior zero is introduced. Choose a closed
interval $J$ about $x_n$, contained in $(1,\infty)$, so small that
\[
\inf_J\Psi_n'>0,
\qquad \sup J<2\inf J.
\]
We restrict $\tau$ further so that $\lambda(\tau)\pm\tau\in J$.
For $s,x\in J$, the ratio $x/s$ is bounded away from $1/2$;
thus $K$ is smooth there. The mean value theorem gives
\begin{equation}
\label{n_dimensional_local_perturbation}
\sup_{x\in J}\left(
 |R_{\tau,\lambda(\tau)}(x)|
 +\left|\frac{d}{dx}R_{\tau,\lambda(\tau)}(x)\right|
 \right)\leq C\tau.
\end{equation}
Here and below the constants may depend on the fixed dimension and
on $J$, but not on sufficiently small $\tau$.
For all $x\geq1$ and $s\in J$,
\[
\frac{\partial K}{\partial s}(s,x)
 =s^{n-1}\bigl(n\Phi_n(x/s)-(x/s)\Phi_n'(x/s)\bigr).
\]
The function $\Phi_n'$ is bounded on $[0,\infty)$ and
$\Phi_n(t)=O(1+t)$, so another application of the mean value theorem
gives the global bound
\begin{equation}
\label{n_dimensional_global_perturbation}
|R_{\tau,\lambda(\tau)}(x)|\leq C\tau(1+x)
                    \qquad(x\geq1).
\end{equation}
By Lemma~\ref{n_dimensional_unique_root} and
\eqref{n_dimensional_psi_infinity}, there is a $c>0$ such that
\[
|\Psi_n(x)|\geq c(1+x)\qquad(x\geq1,\ x\notin J).
\]
For sufficiently small $\tau>0$,
\eqref{n_dimensional_global_perturbation} preserves this nonzero
sign off $J$, while \eqref{n_dimensional_local_perturbation}
ensures that $\Psi_n+R_{\tau,\lambda(\tau)}$ is strictly increasing
on $J$. In view of \eqref{n_dimensional_matching_equation}, we have
proved
\begin{equation}
\label{n_dimensional_only_exterior_zero}
\mathcal D_F(x)=0,\quad x\geq1
       \quad\Longleftrightarrow\quad x=\lambda(\tau).
\end{equation}
Fix such a $\tau$ and write $\lambda=\lambda(\tau)$.

If $x\geq1$ and $x\ne\lambda$, the derivative in
\eqref{n_dimensional_diagonal_derivative} is nonzero, so $r=x$
is not a maximizing radius. At $x=\lambda$, the center is in the
interior of the outer shell. Hence $M_cF(\lambda e_1)=1$.
The ball $B(\lambda e_1,\lambda)$ meets $B_\delta$ in positive
measure, so $A_F(\lambda,\lambda)<1$.
This proves \eqref{n_dimensional_strict_gap} for $x\geq1$.

If $\delta<x<1$, the center is in the interior of $U_\delta$;
again $M_cF(xe_1)=1$, whereas $B(xe_1,x)$ meets the hole in
positive measure. Finally, if $0<x\leq\delta$, then
\[
B\left(\frac{x}{2}e_1,\frac{x}{2}\right)
       \subset B(xe_1,x)\cap B_\delta.
\]
Since $0\leq F\leq1$, this gives $A_F(x,x)\leq1-2^{-n}$.
The ball $B(xe_1,1-x)$ is contained in $B_1$ and contains the
whole hole: indeed, $1-2x\geq1-2\delta>\delta$, because
$\delta<1/3$. Consequently,
\[
A_F(x,1-x)=1-\left(\frac{\delta}{1-x}\right)^n
       >1-2^{-n},
\]
where we used $\delta/(1-x)\leq\delta/(1-\delta)<1/2$.
This proves \eqref{n_dimensional_strict_gap} for all $x>0$.
\end{proof}

\subsection{A uniform exclusion interval and its stability}

The next lemma makes precise the passage from the strict inequality
at $r=x$ to an interval of excluded radii. The perturbation need not
be radial or compactly supported. Its $L^\infty$ norm controls the
averages on bounded ranges of centers, and its $L^1$ norm controls
them at infinity.

\begin{lemma}
\label{n_dimensional_stability}
Suppose that $F$ is radial, $0\leq F\leq1$, and $F$ vanishes on
$B_\delta$ and outside $B_S$, where $0<\delta<S$. Suppose also that
$T=\int_{\mathbb{R}^n}F>0$ and
\[
A_F(x,x)<M_cF(xe_1)\qquad(x>0).
\]
There exist $0<\eta<\delta$ and constants
$\kappa_\infty,\kappa_1>0$ such that, for every nonnegative
measurable function $p$ satisfying
\begin{equation}
\label{n_dimensional_small_perturbation}
\|p\|_\infty\leq\kappa_\infty,
\qquad \|p\|_1\leq\kappa_1,
\end{equation}
we have
\begin{equation}
\label{n_dimensional_safe_radii}
M_c(F+p)(X)=
\sup_{\substack{r>0\\|r-\|X\||>\eta}}
 \frac{1}{c_nr^n}\int_{B(X,r)}(F(y)+p(y))\,dy
       \qquad(X\in\mathbb{R}^n).
\end{equation}
\end{lemma}

\begin{proof}
Write $a=\delta/4$ and choose $R>\max\{1,a\}$ so large that
\begin{equation}
\label{n_dimensional_tail_choice}
\frac34\left(\frac{R+S}{R-1}\right)^n<1.
\end{equation}
For each fixed $r>0$, $A_F(x,r)$ is continuous in $x$; hence
$M_cF(xe_1)$ is lower semicontinuous. Since $A_F(x,x)$ is continuous
for $x>0$, the hypothesis implies
\begin{equation}
\label{n_dimensional_compact_gap}
\gamma=\min_{a\leq x\leq R}
       \bigl(M_cF(xe_1)-A_F(x,x)\bigr)>0.
\end{equation}
Choose $\eta>0$ so small that
\begin{equation}
\label{n_dimensional_gap_width}
\eta\leq\min\left\{
 1,\frac{a}{2},\frac{a\gamma}{8n},
 \frac{T}{16c_{n-1}S^{n-1}}\right\},
\end{equation}
and put
\begin{equation}
\label{n_dimensional_stability_constants}
\kappa_\infty=\min\left\{\frac{\gamma}{4},
                   \frac{T}{2c_n(S+a)^n}\right\},
\qquad \kappa_1=\frac{T}{8}.
\end{equation}
All these constants are strictly positive.
By rotation, it suffices to write $X=xe_1$ with $x\geq0$; rotation
leaves $F$ and the two bounds on $p$ unchanged.

\medskip
\noindent\textit{Centers with $0\leq x\leq a$.}
If $|r-x|\leq\eta$, then
\[
x+r\leq2a+\eta\leq\frac{5\delta}{8}<\delta.
\]
Thus the averaging ball lies in the hole and its average of $F+p$
is at most $\kappa_\infty$. The nonexcluded radius $x+S$ contains
the support of $F$ and gives an average at least
\[
\frac{T}{c_n(x+S)^n}\geq\frac{T}{c_n(a+S)^n}
       >\kappa_\infty.
\]
This excludes the whole interval of radii in this range, including
possible suprema approached as $r\downarrow0$ when $x=0$.

\medskip
\noindent\textit{Centers with $a\leq x\leq R$.}
Since $0\leq F\leq1$, the radial average is locally absolutely
continuous in $r$ and satisfies
\[
\left|\frac{\partial A_F}{\partial r}(x,r)\right|
       \leq\frac{n}{r}
       \qquad\text{for almost every }r>0.
\]
To see this, differentiate the volume-normalized integral: its
derivative is $n/r$ times the difference of the mean of $F$ on
the boundary sphere and its mean on the ball, both lying in $[0,1]$.
For $|r-x|\leq\eta$, every intermediate radius is at least $a/2$.
Consequently, \eqref{n_dimensional_gap_width} gives
\[
|A_F(x,r)-A_F(x,x)|\leq\frac{2n}{a}\eta\leq\frac{\gamma}{4}.
\]
The perturbed average therefore satisfies
\begin{align*}
\frac{1}{c_nr^n}\int_{B(xe_1,r)}(F+p)
 &\leq A_F(x,x)+\frac{\gamma}{4}+\kappa_\infty\\
 &\leq M_cF(xe_1)-\frac{\gamma}{2}.
\end{align*}
For $F$ alone, the excluded averages are at most
$M_cF(xe_1)-3\gamma/4$, so its supremum can already be taken over
nonexcluded radii. Since $p\geq0$, the supremum of the perturbed
averages over those radii is at least $M_cF(xe_1)$. Thus the excluded
radii cannot contribute to the maximal function of $F+p$.

\medskip
\noindent\textit{Centers with $x\geq R$.}
Write $r=x+t$, where $|t|\leq\eta$. For
$y=(y_1,\ldots,y_n)\in B(xe_1,r)$,
\[
y_1>\frac{\|y\|^2-2xt-t^2}{2x}
       \geq-\eta-\frac{\eta^2}{2x}.
\]
Radial symmetry puts exactly half the mass of $F$ in $\{y_1>0\}$.
The additional slab has width at most
$\eta+\eta^2/(2x)\leq2\eta$, and every section of $B_S$ parallel
to $\{y_1=0\}$ has $(n-1)$-dimensional volume at most
$c_{n-1}S^{n-1}$. Since $F\leq1$,
\begin{align*}
\int_{B(xe_1,r)}(F+p)
 &\leq\frac{T}{2}+2c_{n-1}S^{n-1}\eta+\|p\|_1\\
 &\leq\frac{T}{2}+\frac{T}{8}+\frac{T}{8}
       =\frac{3T}{4}.
\end{align*}
As $r\geq x-\eta\geq x-1$, it follows from
\eqref{n_dimensional_tail_choice} that
\[
\frac{1}{c_nr^n}\int_{B(xe_1,r)}(F+p)
 \leq\frac{3T}{4c_n(x-\eta)^n}
 <\frac{T}{c_n(x+S)^n}.
\]
The last expression is the average of $F$ at the nonexcluded radius
$x+S$, and the average of $F+p$ at that radius is no smaller.
This excludes all the remaining radii and proves
\eqref{n_dimensional_safe_radii}.
\end{proof}

\subsection{Completion of the construction}

\begin{proof}[Proof of Theorem~\ref{n_dimensional_theorem}]
Choose $\tau$ and $\lambda$ as in
Lemma~\ref{n_dimensional_matched_shell}. The resulting background
$F$ satisfies the hypotheses of Lemma~\ref{n_dimensional_stability}
with $S=\lambda+\tau$ and
\begin{equation}
\label{n_dimensional_total_mass}
T=c_n\bigl(1-\delta^n+(\lambda+\tau)^n-(\lambda-\tau)^n\bigr).
\end{equation}
Take $\eta$, $\kappa_\infty$, and $\kappa_1$ as in that lemma.
In particular, $\eta\leq\delta/8<\log2$.

The annulus $B_\eta\setminus B_{\eta/2}$ has $2^n-1$ times the
volume of $B_{\eta/2}$. Thus the function
\eqref{n_dimensional_perturbation} satisfies
\begin{equation}
\label{n_dimensional_zero_mass}
\int_{\mathbb{R}^n}h_n=0,
\qquad \|h_n\|_1=c_n\left(\frac{\eta}{2}\right)^n,
\qquad |h_n|\leq\frac12.
\end{equation}
Put $V_\pm(y)=e^{-\|y\|}\pm h_n(y)$. On $B_\eta$ we have
$e^{-\|y\|}\geq e^{-\eta}>1/2$, and outside this ball $h_n=0$.
Consequently,
\[
0<V_\pm\leq\frac32,
\qquad
\|V_\pm\|_1=\int_{\mathbb{R}^n}e^{-\|y\|}\,dy=c_n n!.
\]
Define
\begin{equation}
\label{n_dimensional_epsilon_choice}
\epsilon_0=\min\left\{\frac{2\kappa_\infty}{3},
                         \frac{\kappa_1}{c_n n!}\right\}>0.
\end{equation}
For $0<\epsilon\leq\epsilon_0$, both perturbations
$p_\pm=\epsilon V_\pm$ satisfy
\eqref{n_dimensional_small_perturbation}. Thus
\eqref{n_dimensional_safe_radii} allows us to compute $M_cf$ and
$M_cg$ using only radii with \mbox{$|r-\|X\||>\eta$}.

If $r<\|X\|-\eta$, the ball $B(X,r)$ misses $B_\eta$.
If $r>\|X\|+\eta$, it contains $B_\eta$ completely.
In either case \eqref{n_dimensional_zero_mass} yields
\[
\int_{B(X,r)}h_n=0,\qquad
\int_{B(X,r)}f=\int_{B(X,r)}g.
\]
Taking the suprema over these radii proves the equality in
\eqref{n_dimensional_conclusion} at every $X\in\mathbb{R}^n$.
The formula for $\|f-g\|_1$ follows from $f-g=2\epsilon h_n$.
Strict positivity, radiality, and boundedness have already been
established, and
\[
\|f\|_1=\|g\|_1=T+\epsilon c_n n!<\infty.
\]
This completes the proof.
\end{proof}

Unlike the numerical parameters in the three-dimensional construction,
the parameters of the two-shell background are determined by the
simple-root equation \eqref{n_dimensional_matching_equation} and a
sufficiently small choice of thickness for each fixed $n$. The proof
checks both that this choice creates no additional exterior stationary
diagonal centers and that the final positive perturbations preserve
the excluded interval of radii.

Together with Theorems~\ref{one_dimensional_theorem},
\ref{two_dimensional_theorem}, and \ref{three_dimensional_theorem},
this establishes that the central Hardy--Littlewood maximal operator
is not injective on the nonnegative cone of $L^1(\mathbb{R}^n)$ for
any integer $n\geq1$.


\section*{Acknowledgements and use of AI}

The author acknowledges substantial mathematical assistance from the
large language model \mbox{GPT-6 Astra Pro} (OpenAI). Starting from the author's earlier
one- and two-dimensional constructions and the idea of perturbations
invisible to the central maximal operator, the model proposed the
constructions and proof arguments for dimensions $n\geq3$ presented in
this paper during an extended interaction with the author. Its
contribution went beyond language editing and included the development
of the radial constructions, the analysis of the relevant ball averages,
and the extension to higher dimensions. The model also assisted with
the detailed exposition of the two-dimensional argument, calculations,
and the preparation and editing of the manuscript.

The author has independently checked all arguments and takes full
responsibility for the final content of the paper.

\begin{thebibliography}{3}

\bibitem{deguzman1975}
Miguel de Guzm\'an, \emph{Differentiation of integrals in $\mathbb{R}^n$},
Lecture Notes in Math., vol.~481, Springer-Verlag, Berlin, 1975.

\bibitem{ephremidze2010}
Lasha Ephremidze, \emph{On the uniqueness property of various maximal operators},
RIMS K\^oky\^uroku Bessatsu \textbf{B22} (2010), 137--144.

\bibitem{hagelsteinparissis2015}
Paul A. Hagelstein and Ioannis Parissis,
\emph{Solyanik estimates and local H\"older continuity of halo functions
of geometric maximal operators},
Adv. Math. \textbf{285} (2015), 434--453.

\end{thebibliography}


\end{document}


